\documentclass[11pt]{article}
\usepackage[utf8]{inputenc}
\usepackage[T1]{fontenc}
\usepackage[spanish,es-noquoting]{babel}
\usepackage{amsmath,amssymb}
\usepackage{newtxtext,newtxmath}
\usepackage{booktabs,array,enumitem}
\usepackage[margin=2.5cm]{geometry}
\usepackage{microtype}
\newif\ifsolucion
\ifdefined\ENUNCIADO\solucionfalse\else\soluciontrue\fi
\setlength{\parskip}{0.4em}
\begin{document}

%============================ ENCABEZADO ============================
\begin{center}
{\normalsize\textbf{(IIP314W-2) [2026-T2] OPTIMIZACIÓN APLICADA A NEGOCIOS}}\\[0.4em]
\rule{\linewidth}{0.8pt}\\[0.6em]
{\Large\bfseries Ayudantía 5: Modelamiento con variables binarias%
\ifsolucion\ --- Pauta (Solución)\fi}\\[0.8em]
\rule{\linewidth}{0.4pt}\\[0.5em]
{\normalsize
Profesor: Rodrigo Trigo Vilches \quad$\cdot$\quad Ayudante: Vicente Ramírez\\[0.2em]
Fecha: 8 de julio de 2026 \quad$\cdot$\quad Universidad del Desarrollo}
\end{center}

\vspace{0.6em}

%============================ OBJETIVOS ============================
\section*{Objetivos de la ayudantía}

Esta ayudantía es de \textbf{modelamiento puro}: no escribiremos código (a diferencia de la Ayudantía 4, que resolvía en Python/Gurobi). El foco está en \emph{traducir} enunciados de negocio a modelos matemáticos con variables binarias, de forma \textbf{algebraica y paramétrica} (sin números concretos). Al terminar, deberías ser capaz de:

\begin{itemize}[leftmargin=1.6em,itemsep=0.2em]
    \item Traducir problemas de negocio a modelos de optimización con \textbf{variables binarias}.
    \item Modelar decisiones de \textbf{sí/no} (abrir una instalación, cargar un módulo, encender una máquina) con una binaria por decisión.
    \item Usar \textbf{restricciones de activación} que ligan una variable binaria con una variable continua (``solo puede fluir/producir si está encendida'').
    \item Aplicar la técnica de \textbf{big-M} y entender por qué conviene elegir la cota más ajustada posible.
    \item Modelar \textbf{costos fijos} (fixed-charge), \textbf{incompatibilidades} (exclusión mutua), \textbf{prerrequisitos} (implicaciones if-then) y \textbf{coberturas} (mínimos por categoría).
    \item Formular todo de forma \textbf{algebraica y paramétrica}, definiendo conjuntos, parámetros, variables, objetivo y restricciones.
\end{itemize}

%============================ REPASO ============================
\section*{Repaso: modelamiento con variables binarias}

\noindent\textbf{1. La variable binaria.} Una variable $y\in\{0,1\}$ representa una decisión \emph{sí/no} indivisible: $y=1$ si se toma la acción (abrir, cargar, encender, seleccionar) y $y=0$ si no. Toda la potencia del modelamiento entero está en \emph{ligar} estas binarias entre sí y con las variables continuas mediante restricciones lineales.

\smallskip
\noindent\textbf{2. Costo fijo (fixed-charge).} Muchas decisiones tienen un costo $f$ por el solo hecho de ``activarse'', independiente del nivel de uso $x$. La estructura canónica combina una binaria de activación $y$ con una variable continua $x$ ligadas por un big-M:
\[
\min \; \dots + f\,y + (\text{costo variable de } x)
\qquad\text{s.a.}\qquad
x \le M\,y,\quad x\ge 0,\quad y\in\{0,1\}.
\]
Si $y=0$, la restricción $x\le M\cdot 0=0$ obliga a $x=0$ (no se usa la opción y no se paga $f$); si $y=1$, se libera $x\le M$ y se paga el costo fijo $f$. Así el modelo ``cobra'' $f$ solo cuando efectivamente se activa la opción.

\smallskip
\noindent\textbf{3. Restricción de activación / big-M.} En general, la frase ``$x$ solo puede ser positiva si $y=1$'' se escribe $x \le M\,y$, con $M$ una cota superior válida de $x$. Cuando además se exige un \emph{mínimo} si la opción está activa (por ejemplo un mínimo técnico), se usa una \textbf{activación doble}:
\[
\ell\,y \;\le\; x \;\le\; u\,y,
\]
que fuerza $x=0$ cuando $y=0$ y $\ell \le x \le u$ cuando $y=1$.

\smallskip
\noindent\textbf{4. Lógica proposicional con binarias.} Las relaciones lógicas entre decisiones se vuelven desigualdades lineales:
\[
\underbrace{y_i \le y_j}_{\text{si }i\text{ entonces }j}
\qquad
\underbrace{y_i + y_j \le 1}_{\text{incompatibles (a lo más uno)}}
\qquad
\underbrace{\textstyle\sum_i y_i \ge 1}_{\text{al menos uno}}
\qquad
\underbrace{\textstyle\sum_i y_i \le k}_{\text{a lo más }k}.
\]
La implicación $y_i\le y_j$ dice ``si se elige $i$ ($y_i=1$) entonces obligatoriamente $y_j=1$''; es de un solo sentido. La exclusión $y_i+y_j\le1$ prohíbe el caso $(1,1)$.

\smallskip
\noindent\textbf{5. Cómo elegir $M$.} El big-M debe ser \emph{lo más pequeño posible pero válido} (una cota superior legítima de la variable que acota). Un $M$ innecesariamente gigante no cambia el óptimo entero, pero \textbf{debilita la relajación lineal} y hace mucho más lento el \emph{branch and bound}. Siempre conviene usar la cota natural del problema (una capacidad, una demanda acumulada, etc.).

%============================ EJERCICIO 1 ============================
\newpage

\section{Ejercicio 1 — El puente aéreo humanitario}

Una organización no gubernamental (ONG) de ayuda de emergencia coordina un \textbf{puente aéreo humanitario} hacia una región afectada por un desastre natural. La ventana de operación es breve: solo se dispone de \textbf{un vuelo} de un avión de carga, y la logística en tierra impide reprogramar despachos. Por eso, la decisión de \emph{qué cargar} en ese único vuelo es crítica y debe tomarse de una sola vez.

La ayuda está preparada en \textbf{módulos} estandarizados listos para despachar. Cada módulo $i$ del conjunto $\mathcal{I}$ es una unidad indivisible —un pallet sellado— que se carga entero o no se carga: no tiene sentido enviar ``medio módulo''. A cada módulo el equipo humanitario le asignó, a partir de la evaluación de necesidades en terreno, un \textbf{impacto humanitario} $v_i$, un puntaje que resume cuántas personas beneficia y con qué urgencia. Además, cada módulo ocupa un \textbf{volumen} $a_i$ (en $\text{m}^3$) en la bodega, tiene una \textbf{masa} $m_i$ (en kg) y significó a la ONG un \textbf{costo de adquisición} $c_i$ para comprarlo y prepararlo.

La ONG quiere elegir el subconjunto de módulos que \textbf{maximice el impacto humanitario total} del vuelo, respetando todas las limitaciones físicas, presupuestarias y operativas. Estas limitaciones surgen de forma natural del problema:

\begin{itemize}[leftmargin=1.4em]
    \item El avión tiene una \textbf{bodega} de capacidad finita: la suma de los volúmenes de lo cargado no puede exceder $V$ metros cúbicos.
    \item El avión tiene un \textbf{payload máximo} (carga útil que puede levantar con seguridad): la masa total no puede superar $M$ kilogramos.
    \item La operación se financia con un \textbf{presupuesto} acotado $B$: la ONG no puede gastar en adquisición más de lo que tiene.
    \item El operativo debe ser \textbf{equilibrado}: los módulos se clasifican en categorías $g$ (agua, alimentos, medicina, refugio y energía), y por criterio humanitario se exige enviar \textbf{al menos} $n_g$ módulos de cada categoría, para no llegar con un avión lleno de agua pero sin medicinas.
    \item Algunos módulos son \textbf{incompatibles} entre sí y no pueden viajar juntos (por ejemplo, un módulo de combustible inflamable junto a uno de material sensible): para ciertos pares $(i,j)\in\mathcal{P}$ se puede cargar a lo más uno de los dos.
    \item Hay \textbf{prerrequisitos} técnicos: ciertos módulos de \emph{medicina refrigerada} solo sirven si viaja también el \textbf{generador} que mantiene la cadena de frío; si se carga el módulo $i$, entonces obligatoriamente debe cargarse el módulo $j$ asociado. Estas dependencias se listan en el conjunto $\mathcal{R}$ de pares $(i,j)$.
    \item Por límites de \textbf{manejo y estiba en tierra} —cuadrillas y tiempo de carga acotados— no se pueden despachar más de $N$ módulos en total, aun si físicamente cupieran.
\end{itemize}

\bigskip
\noindent\textbf{Conjuntos}

\begin{tabular}{@{}ll@{}}
\toprule
Símbolo & Descripción \\
\midrule
$\mathcal{I}$ & Conjunto de módulos de ayuda disponibles, indexados por $i$ \\
$\mathcal{G}$ & Conjunto de categorías de ayuda (agua, alimentos, medicina, refugio, energía), indexadas por $g$ \\
$\mathcal{I}_g \subseteq \mathcal{I}$ & Subconjunto de módulos que pertenecen a la categoría $g$ \\
$\mathcal{P}$ & Conjunto de pares $(i,j)$ de módulos mutuamente incompatibles \\
$\mathcal{R}$ & Conjunto de pares $(i,j)$ tales que cargar $i$ requiere cargar $j$ (prerrequisito) \\
\bottomrule
\end{tabular}

\bigskip
\noindent\textbf{Parámetros}

\begin{tabular}{@{}ll@{}}
\toprule
Símbolo & Descripción \\
\midrule
$v_i$ & Impacto humanitario del módulo $i$ \\
$a_i$ & Volumen que ocupa el módulo $i$ $[\text{m}^3]$ \\
$m_i$ & Masa del módulo $i$ $[\text{kg}]$ \\
$c_i$ & Costo de adquisición del módulo $i$ $[\$]$ \\
$V$ & Capacidad de volumen de la bodega del avión $[\text{m}^3]$ \\
$M$ & Payload máximo (masa) que puede transportar el avión $[\text{kg}]$ \\
$B$ & Presupuesto total de adquisición $[\$]$ \\
$n_g$ & Número mínimo de módulos exigido de la categoría $g$ \\
$N$ & Número máximo de módulos que pueden despacharse en el vuelo \\
\bottomrule
\end{tabular}

\ifsolucion

\subsection*{Variables de decisión}

La decisión es, para cada módulo, \emph{cargarlo o no cargarlo}: un ``sí/no'' indivisible. Esto se modela con una variable \textbf{binaria} por módulo:
\[
y_i =
\begin{cases}
1 & \text{si el módulo } i \text{ se carga en el vuelo,}\\[2pt]
0 & \text{en caso contrario,}
\end{cases}
\qquad \forall i \in \mathcal{I}.
\]

\subsection*{Función objetivo}

Se busca maximizar el impacto humanitario total del cargamento, esto es, la suma de los impactos de los módulos efectivamente cargados:
\[
\max \; \sum_{i \in \mathcal{I}} v_i \, y_i .
\]

\subsection*{Restricciones}

\begin{align}
\sum_{i \in \mathcal{I}} a_i \, y_i &\le V
    & &\text{(volumen)} \label{eq:vol}\\[4pt]
\sum_{i \in \mathcal{I}} m_i \, y_i &\le M
    & &\text{(masa)} \label{eq:masa}\\[4pt]
\sum_{i \in \mathcal{I}} c_i \, y_i &\le B
    & &\text{(presupuesto)} \label{eq:pres}\\[4pt]
\sum_{i \in \mathcal{I}_g} y_i &\ge n_g
    & &\text{(cobertura)} \quad \forall g \in \mathcal{G} \label{eq:cob}\\[4pt]
y_i + y_j &\le 1
    & &\text{(incompatibilidad)} \quad \forall (i,j) \in \mathcal{P} \label{eq:incomp}\\[4pt]
y_i &\le y_j
    & &\text{(prerrequisito)} \quad \forall (i,j) \in \mathcal{R} \label{eq:prereq}\\[4pt]
\sum_{i \in \mathcal{I}} y_i &\le N
    & &\text{(cardinalidad)} \label{eq:card}\\[4pt]
y_i &\in \{0,1\}
    & &\text{(naturaleza)} \quad \forall i \in \mathcal{I} \label{eq:bin}
\end{align}

\subsection*{Comentario de modelamiento}

Este es un problema de \textbf{mochila (knapsack) multidimensional con lógica de selección}: la mochila es el avión, y ``multidimensional'' porque hay tres presupuestos simultáneos que respetar —volumen \eqref{eq:vol}, masa \eqref{eq:masa} y dinero \eqref{eq:pres}—, todos con la misma estructura $\sum(\text{recurso})\,y_i \le (\text{disponible})$. Lo interesante no es la capacidad, que ya vimos como restricción clásica, sino los tres \emph{trucos de modelamiento} con binarias que aparecen:

\medskip
\noindent\textbf{Incompatibilidad (mutuamente excluyentes).} Cuando dos módulos $i,j$ no pueden viajar juntos, imponemos $y_i + y_j \le 1$ en \eqref{eq:incomp}. La lectura es directa: la suma de dos variables $0$--$1$ solo alcanza el valor $2$ cuando ambas valen $1$; al acotarla por $1$ prohibimos exactamente ese caso, dejando libres las tres combinaciones admisibles $(0,0)$, $(1,0)$ y $(0,1)$. Es la forma canónica de decir ``a lo más uno de los dos''.

\medskip
\noindent\textbf{Prerrequisito o implicación if-then.} La frase ``si cargo la medicina refrigerada $i$, \emph{entonces} debo cargar el generador $j$'' es una implicación lógica $y_i = 1 \Rightarrow y_j = 1$. Se traduce en la desigualdad $y_i \le y_j$ de \eqref{eq:prereq}. Conviene verificar los cuatro casos: si $y_i=1$, la restricción fuerza $y_j \ge 1$, es decir $y_j=1$ (el generador viaja sí o sí); si $y_i=0$, entonces $y_j$ queda libre (puedo llevar el generador aunque no lleve esa medicina, o no llevar ninguno). Nótese que la implicación es \emph{en un solo sentido}: llevar el generador no obliga a llevar la medicina.

\medskip
\noindent\textbf{Cobertura por categoría.} La exigencia de equilibrio humanitario se modela con \eqref{eq:cob}, una restricción $\ge$ por cada categoría $g$, que suma las binarias \emph{solo sobre los módulos de esa categoría} (el subconjunto $\mathcal{I}_g$) y obliga a que haya al menos $n_g$ de ellos. Aquí el papel de los subconjuntos $\mathcal{I}_g$ es clave: nos permite escribir una familia compacta de restricciones —una por categoría— sin enumerar módulo por módulo. Finalmente, la cardinalidad \eqref{eq:card} es la ``otra cara'' de la cobertura: un techo global $\le N$ sobre \emph{todas} las binarias que limita cuántos módulos entran en total, capturando la restricción operativa de estiba que no es física (volumen o masa) sino de manejo en tierra.

\fi

%============================ EJERCICIO 2 ============================
\newpage

\section{Ejercicio 2 --- Rediseño de una red logística en el tiempo}

Una empresa nacional de distribución de bienes de consumo enfrenta una decisión estratégica que rara vez se toma mirando un solo mes: quiere \textbf{rediseñar toda su red logística} a lo largo de un horizonte de planificación de varios periodos. La red tiene tres eslabones. En el primero están las \textbf{plantas} productoras $\mathcal{P}$, que son \emph{candidatas} --- algunas ya operan, otras podrían construirse ---; en el segundo, las \textbf{bodegas} $\mathcal{B}$, centros de acopio regionales que reciben la producción y la almacenan; y en el tercero, las \textbf{zonas de demanda} $\mathcal{J}$, los mercados que hay que abastecer. El producto fluye siempre en el mismo sentido: planta $\to$ bodega $\to$ zona.

Lo que hace interesante (y difícil) al problema es que la empresa no decide una foto fija, sino una \emph{película}: en cada periodo $t\in\mathcal{T}$ puede tener ciertas plantas y bodegas \textbf{operativas} y otras cerradas, y puede \textbf{abrir} instalaciones que estaban cerradas o \textbf{cerrar} las que estaban abiertas, pagando cada vez el costo de la transición. El mercado tampoco es estático: cada zona $j$ exige una demanda $d_{jt}$ que cambia periodo a periodo (estacionalidad, crecimiento, campañas).

Cada instalación operativa cuesta dinero solo por estar encendida: mantener una planta $p$ abierta durante el periodo $t$ cuesta $\mathit{FP}_{pt}$, y una bodega $b$ cuesta $\mathit{FB}_{bt}$. Encender por primera vez (o reabrir) una planta tiene un costo de apertura $\mathit{AP}_{p}$; apagarla cuesta $\mathit{CP}_{p}$ (desmovilizar personal, indemnizaciones, mantención en frío); las bodegas tienen sus equivalentes $\mathit{AB}_{b}$ y $\mathit{CB}_{b}$. Del lado operativo, cada planta operativa produce a lo más $\mathit{Cap}_{pt}$ unidades en el periodo $t$; cada bodega puede almacenar a lo más $\mathit{Alm}_{b}$ unidades; enviar una unidad de la planta $p$ a la bodega $b$ cuesta $\mathit{cp}_{pb}$ (producción más transporte del primer tramo) y de la bodega $b$ a la zona $j$ cuesta $\mathit{cg}_{bj}$. Finalmente, lo que una bodega no despacha en un periodo queda como \textbf{inventario}, que se arrastra al periodo siguiente pagando un costo de mantención $h_{b}$ por unidad almacenada al cierre del periodo.

La gerencia quiere el plan de \emph{apertura, cierre, producción, transporte y almacenamiento} que satisface toda la demanda a lo largo del horizonte al \textbf{mínimo costo total}.

\medskip

\begin{center}
\begin{tabular}{@{}ll@{}}
\toprule
\textbf{Conjunto} & \textbf{Descripción} \\
\midrule
$\mathcal{P}$ & plantas candidatas (índice $p$) \\
$\mathcal{B}$ & bodegas candidatas (índice $b$) \\
$\mathcal{J}$ & zonas de demanda (índice $j$) \\
$\mathcal{T}$ & periodos del horizonte, $t=1,\dots,|\mathcal{T}|$ (índice $t$) \\
\bottomrule
\end{tabular}
\end{center}

\medskip

\begin{center}
\begin{tabular}{@{}ll@{}}
\toprule
\textbf{Parámetro} & \textbf{Significado} \\
\midrule
$d_{jt}$ & demanda de la zona $j$ en el periodo $t$ $[\text{unid}]$ \\
$\mathit{FP}_{pt}$ & costo fijo de operar la planta $p$ en el periodo $t$ $[\$]$ \\
$\mathit{FB}_{bt}$ & costo fijo de operar la bodega $b$ en el periodo $t$ $[\$]$ \\
$\mathit{AP}_{p},\ \mathit{CP}_{p}$ & costo de abrir / cerrar la planta $p$ $[\$]$ \\
$\mathit{AB}_{b},\ \mathit{CB}_{b}$ & costo de abrir / cerrar la bodega $b$ $[\$]$ \\
$\mathit{Cap}_{pt}$ & capacidad de producción de la planta $p$ en $t$ $[\text{unid}]$ \\
$\mathit{Alm}_{b}$ & capacidad de almacenamiento de la bodega $b$ $[\text{unid}]$ \\
$\mathit{cp}_{pb}$ & costo unitario producción $+$ transporte planta $p\to$ bodega $b$ $[\$/\text{unid}]$ \\
$\mathit{cg}_{bj}$ & costo unitario de transporte bodega $b\to$ zona $j$ $[\$/\text{unid}]$ \\
$h_{b}$ & costo de mantener una unidad de inventario en la bodega $b$ $[\$/\text{unid}]$ \\
$y_{p0},\ z_{b0}$ & estado inicial (dado) de planta $p$ / bodega $b$ antes del horizonte \\
$s_{b0}$ & inventario inicial de la bodega $b$ (dado) $[\text{unid}]$ \\
$M^{P}_{pt},\ M^{B}_{bt}$ & cotas grandes (big-M) para ligar flujo con apertura \\
\bottomrule
\end{tabular}
\end{center}

\ifsolucion

\subsection*{Variables de decisión}

\noindent\textbf{Binarias de estado y de transición.}
\begin{align*}
y_{pt} &=
\begin{cases}
1 & \text{si la planta } p \text{ está \emph{operativa} en el periodo } t,\\
0 & \text{en caso contrario},
\end{cases}
& \forall p\in\mathcal{P},\ t\in\mathcal{T},\\[2pt]
z_{bt} &=
\begin{cases}
1 & \text{si la bodega } b \text{ está \emph{operativa} en el periodo } t,\\
0 & \text{en caso contrario},
\end{cases}
& \forall b\in\mathcal{B},\ t\in\mathcal{T},\\[2pt]
o^{P}_{pt} &= 1 \text{ si la planta } p \text{ se \emph{abre} al inicio de } t;\quad
c^{P}_{pt} = 1 \text{ si se \emph{cierra}}, & \forall p\in\mathcal{P},\ t\in\mathcal{T},\\[2pt]
o^{B}_{bt} &= 1 \text{ si la bodega } b \text{ se \emph{abre} al inicio de } t;\quad
c^{B}_{bt} = 1 \text{ si se \emph{cierra}}, & \forall b\in\mathcal{B},\ t\in\mathcal{T}.
\end{align*}

\noindent\textbf{Continuas no negativas (flujos e inventario).}
\begin{align*}
f_{pbt} &\ge 0:\ \text{unidades enviadas de la planta } p \text{ a la bodega } b \text{ en } t, & \forall p,b,t,\\
g_{bjt} &\ge 0:\ \text{unidades enviadas de la bodega } b \text{ a la zona } j \text{ en } t, & \forall b,j,t,\\
s_{bt} &\ge 0:\ \text{inventario al final del periodo } t \text{ en la bodega } b, & \forall b,t.
\end{align*}

\subsection*{Función objetivo}

Se minimiza el costo total del horizonte, compuesto por cinco bloques: (i) costos fijos de operación de plantas y bodegas; (ii) costos de apertura; (iii) costos de cierre; (iv) costos de producción y transporte en los dos tramos; y (v) costo de mantener inventario.

\begin{align*}
\min \quad
&\sum_{t\in\mathcal{T}}\sum_{p\in\mathcal{P}} \mathit{FP}_{pt}\,y_{pt}
+ \sum_{t\in\mathcal{T}}\sum_{b\in\mathcal{B}} \mathit{FB}_{bt}\,z_{bt}
&& \text{(operación)}\\
&+ \sum_{t\in\mathcal{T}}\sum_{p\in\mathcal{P}}\!\big(\mathit{AP}_{p}\,o^{P}_{pt} + \mathit{CP}_{p}\,c^{P}_{pt}\big)
+ \sum_{t\in\mathcal{T}}\sum_{b\in\mathcal{B}}\!\big(\mathit{AB}_{b}\,o^{B}_{bt} + \mathit{CB}_{b}\,c^{B}_{bt}\big)
&& \text{(transiciones)}\\
&+ \sum_{t\in\mathcal{T}}\sum_{p\in\mathcal{P}}\sum_{b\in\mathcal{B}} \mathit{cp}_{pb}\,f_{pbt}
+ \sum_{t\in\mathcal{T}}\sum_{b\in\mathcal{B}}\sum_{j\in\mathcal{J}} \mathit{cg}_{bj}\,g_{bjt}
&& \text{(prod.\ + transporte)}\\
&+ \sum_{t\in\mathcal{T}}\sum_{b\in\mathcal{B}} h_{b}\,s_{bt}.
&& \text{(inventario)}
\end{align*}

\subsection*{Restricciones}

\begin{align}
\sum_{b\in\mathcal{B}} f_{pbt} &\le \mathit{Cap}_{pt}\,y_{pt}
&& \forall p\in\mathcal{P},\ t\in\mathcal{T}
&& \text{(cap.\ planta)} \label{c:capplanta}\\[2pt]
s_{bt} &\le \mathit{Alm}_{b}\,z_{bt}
&& \forall b\in\mathcal{B},\ t\in\mathcal{T}
&& \text{(cap.\ bodega)} \label{c:capbodega}\\[2pt]
\sum_{p\in\mathcal{P}} f_{pbt} &\le M^{B}_{bt}\,z_{bt}
&& \forall b\in\mathcal{B},\ t\in\mathcal{T}
&& \text{(activa bodega)} \label{c:activabodega}\\[2pt]
s_{bt} &= s_{b,t-1} + \sum_{p\in\mathcal{P}} f_{pbt} - \sum_{j\in\mathcal{J}} g_{bjt}
&& \forall b\in\mathcal{B},\ t\in\mathcal{T}
&& \text{(balance inv.)} \label{c:balance}\\[2pt]
\sum_{b\in\mathcal{B}} g_{bjt} &\ge d_{jt}
&& \forall j\in\mathcal{J},\ t\in\mathcal{T}
&& \text{(demanda)} \label{c:demanda}\\[2pt]
y_{pt} - y_{p,t-1} &= o^{P}_{pt} - c^{P}_{pt}
&& \forall p\in\mathcal{P},\ t\in\mathcal{T}
&& \text{(trans.\ planta)} \label{c:transp}\\[2pt]
o^{P}_{pt} + c^{P}_{pt} &\le 1
&& \forall p\in\mathcal{P},\ t\in\mathcal{T}
&& \text{(exclus.\ planta)} \label{c:exclp}\\[2pt]
z_{bt} - z_{b,t-1} &= o^{B}_{bt} - c^{B}_{bt}
&& \forall b\in\mathcal{B},\ t\in\mathcal{T}
&& \text{(trans.\ bodega)} \label{c:transb}\\[2pt]
o^{B}_{bt} + c^{B}_{bt} &\le 1
&& \forall b\in\mathcal{B},\ t\in\mathcal{T}
&& \text{(exclus.\ bodega)} \label{c:exclb}\\[2pt]
y_{pt},\, z_{bt},\, o^{P}_{pt},\, c^{P}_{pt},\, o^{B}_{bt},\, c^{B}_{bt} &\in \{0,1\}
&& \forall p,b,t
&& \text{(binarias)} \label{c:bin}\\[2pt]
f_{pbt},\, g_{bjt},\, s_{bt} &\ge 0
&& \forall p,b,j,t
&& \text{(no neg.)} \label{c:noneg}
\end{align}

\noindent En las transiciones \eqref{c:transp} y \eqref{c:transb}, el término $y_{p,0}$ (resp.\ $z_{b,0}$) es el \emph{estado inicial dado} de la instalación; y en el balance \eqref{c:balance}, $s_{b,0}$ es el \emph{inventario inicial dado}. Ambos son parámetros, no variables.

\subsection*{Comentario de modelamiento}

Este modelo combina las tres ideas que hacen difícil (y realista) a la planificación de redes, y vale la pena aislar cada una.

\textbf{El big-M que liga el flujo con la binaria de apertura.} Una planta no puede despachar nada si no está encendida. La restricción \eqref{c:capplanta} logra exactamente eso de manera elegante: cuando $y_{pt}=0$, el lado derecho vale $\mathit{Cap}_{pt}\cdot 0 = 0$, y como los flujos son no negativos, se fuerza $\sum_b f_{pbt}=0$; cuando $y_{pt}=1$, el lado derecho recupera la capacidad real $\mathit{Cap}_{pt}$ y la restricción cumple su segundo rol, el de \emph{cota de producción}. Es decir, la misma desigualdad hace dos trabajos: \emph{activar} (apagar el flujo si está cerrada) y \emph{limitar} (respetar la capacidad si está abierta), usando la propia capacidad como su big-M natural. En la bodega separamos ambos roles: \eqref{c:capbodega} liga el \emph{inventario} con la apertura (una bodega cerrada no guarda stock) y \eqref{c:activabodega} usa un big-M $M^{B}_{bt}$ para apagar el \emph{flujo entrante} si la bodega no opera. Un valor razonable de este big-M es $M^{B}_{bt}=\min\{\mathit{Alm}_b,\ \sum_j d_{jt'}\ \text{acumulada}\}$ o, más simple, la suma de capacidades de las plantas; conviene tomarlo lo más ajustado posible, porque un big-M innecesariamente grande relaja la formulación y hace más lento el \emph{branch and bound}.

\textbf{La lógica de transición abrir/cerrar entre periodos.} La pareja \eqref{c:transp}--\eqref{c:exclp} es el corazón temporal del modelo. La ecuación $y_{pt}-y_{p,t-1}=o^{P}_{pt}-c^{P}_{pt}$ traduce el cambio de estado en eventos: si la planta pasa de apagada a encendida ($0\to1$) la diferencia es $+1$ y obliga a $o^{P}_{pt}=1$ (hubo apertura); si pasa de encendida a apagada ($1\to0$) la diferencia es $-1$ y obliga a $c^{P}_{pt}=1$ (hubo cierre); si no cambia, ambas quedan en cero. La desigualdad de exclusividad $o^{P}_{pt}+c^{P}_{pt}\le 1$ impide la solución absurda de ``abrir y cerrar simultáneamente'' para inflar costos o burlar el balance. Gracias a esta contabilidad, la función objetivo cobra $\mathit{AP}_p$ y $\mathit{CP}_p$ \emph{solo en el periodo exacto} en que ocurre la transición, capturando que reconfigurar la red no es gratis y desincentivando el ``prende y apaga'' oportunista.

\textbf{El balance de inventario multiperiodo.} La restricción \eqref{c:balance} es lo que convierte un conjunto de problemas de transporte independientes en un \emph{plan} acoplado en el tiempo: el inventario al cierre de $t$ es el que había al cierre de $t-1$, más lo que entró desde las plantas, menos lo que salió hacia las zonas. Este arrastre permite que una bodega \emph{acumule} producción barata en periodos de baja demanda para volcarla cuando la demanda sube, siempre pagando $h_b$ por cada unidad guardada. Es la misma estructura de conservación de flujo que en un problema de lote económico, pero aquí anidada dentro de decisiones de apertura: solo una bodega abierta (por \eqref{c:capbodega}) puede sostener inventario, y solo una planta abierta (por \eqref{c:capplanta}) puede alimentarla.

Finalmente, la restricción de \textbf{demanda} \eqref{c:demanda} es la exigencia de servicio del negocio (se puede escribir con $=$ si no se permite sobre-servir), y cierra la cadena obligando a que \emph{alguna} bodega abierta abastezca cada zona en cada periodo. El modelo así integra las decisiones estratégicas (qué construir y cuándo), tácticas (cuánto almacenar) y operativas (cómo transportar) en una sola optimización coherente.

\fi

%============================ EJERCICIO 3 ============================
\newpage

\section{Ejercicio 3 — Encendido y despacho de un parque de generadoras}

Un \textbf{operador eléctrico} debe programar, al \textbf{mínimo costo}, el funcionamiento de su parque de centrales de generación a lo largo del día siguiente. El día se divide en un conjunto de \textbf{periodos} $\mathcal{T}$ (bloques horarios) y, en cada periodo $t$, el operador conoce con antelación la \textbf{demanda} de energía $d_t$ que el sistema deberá abastecer sin falla. Para satisfacerla dispone de un conjunto de \textbf{generadoras} $\mathcal{G}$ —centrales térmicas, de gas y de respaldo—, cada una con características técnicas y económicas propias.

La particularidad de este problema, y lo que lo distingue de un despacho puramente continuo, es que \textbf{encender una central no es gratis ni instantáneo}. Toda generadora tiene un \textbf{mínimo técnico} $\mathit{Pmin}_g$: si está funcionando, por razones físicas (estabilidad de la turbina, temperatura de la caldera) no puede entregar menos que ese piso; y tiene una \textbf{potencia máxima} $\mathit{Pmax}_g$ que no puede superar. Además, mantener una central encendida cuesta $\mathit{cf}_g$ por periodo aunque casi no genere (el llamado costo de \emph{no-load}: personal, combustible de sostenimiento, servicios auxiliares), la energía efectivamente producida cuesta $\mathit{cv}_g$ por unidad, y cada vez que una central pasa de apagada a encendida se incurre en un costo de \textbf{arranque} $\mathit{ca}_g$ (calentar la caldera, sincronizar con la red), que puede ser muy elevado.

El operador enfrenta, entonces, un dilema económico clásico: apagar una central cara en las horas de baja demanda ahorra su costo fijo y variable, pero obliga a pagar de nuevo el arranque cuando la demanda vuelva a subir. Decidir \emph{cuándo} encender y apagar cada unidad —y \emph{cuánto} generar en cada periodo— es el problema de \textbf{unit commitment} (compromiso de unidades). A esto se suman dos exigencias operativas de seguridad: por confiabilidad, la capacidad total \emph{encendida} en cada periodo debe superar la demanda en un margen de \textbf{reserva} $\mathit{rr}_t$ (para absorber fallas o peaks imprevistos), y cada central tiene una \textbf{rampa} $R_g$ que limita cuánto puede variar su potencia de un periodo al siguiente (una turbina no sube ni baja carga de golpe).

Formule un modelo de optimización que le diga al operador, para cada generadora y cada periodo, si debe estar encendida, cuándo arranca y cuánta potencia entrega, minimizando el costo total del día.

\begin{table}[h]
\centering
\begin{tabular}{@{}ll@{}}
\toprule
\textbf{Conjunto} & \textbf{Descripción} \\
\midrule
$\mathcal{G}$ & generadoras del parque (índice $g$) \\
$\mathcal{T}$ & periodos del horizonte de planificación (índice $t$) \\
\bottomrule
\end{tabular}
\end{table}

\begin{table}[h]
\centering
\begin{tabular}{@{}ll@{}}
\toprule
\textbf{Parámetro} & \textbf{Significado} \\
\midrule
$d_t$ & demanda de energía a abastecer en el periodo $t$ $[\text{MW}]$ \\
$\mathit{Pmin}_g$ & mínimo técnico de la generadora $g$ (potencia mínima si está encendida) $[\text{MW}]$ \\
$\mathit{Pmax}_g$ & potencia máxima de la generadora $g$ $[\text{MW}]$ \\
$\mathit{cv}_g$ & costo variable de $g$ por unidad de energía generada $[\$/\text{MW}]$ \\
$\mathit{cf}_g$ & costo fijo de $g$ por periodo estando encendida (no-load) $[\$]$ \\
$\mathit{ca}_g$ & costo de arranque de $g$ cada vez que se enciende $[\$]$ \\
$\mathit{rr}_t$ & reserva operativa exigida en el periodo $t$ $[\text{MW}]$ \\
$R_g$ & rampa: variación máxima de potencia de $g$ entre periodos consecutivos $[\text{MW}]$ \\
$u^0_g$ & estado inicial de $g$ (1 si venía encendida del día anterior, 0 si no) \\
$p^0_g$ & potencia inicial de $g$ al comienzo del horizonte $[\text{MW}]$ \\
\bottomrule
\end{tabular}
\end{table}

\ifsolucion

\subsection*{Variables de decisión}

Para capturar simultáneamente el \emph{estado} de cada central, su \emph{transición} y su \emph{nivel de generación}, el modelo combina variables binarias con variables continuas:

\begin{align*}
u_{gt} &\in \{0,1\}: & &\text{$1$ si la generadora $g$ está \textbf{encendida} en el periodo $t$, $0$ si no;} \\
v_{gt} &\in \{0,1\}: & &\text{$1$ si la generadora $g$ \textbf{arranca} (pasa de apagada a encendida) al inicio de $t$;} \\
p_{gt} &\ge 0: & &\text{\textbf{potencia} entregada por la generadora $g$ durante el periodo $t$ $[\text{MW}]$.}
\end{align*}

\subsection*{Función objetivo}

Se minimiza el costo total del horizonte, suma de tres componentes: el costo variable de la energía efectivamente generada, el costo fijo de mantener encendidas las unidades y el costo de los arranques.

\begin{align*}
\min \; \sum_{g\in\mathcal{G}}\sum_{t\in\mathcal{T}} \Big( \mathit{cv}_g\, p_{gt} \;+\; \mathit{cf}_g\, u_{gt} \;+\; \mathit{ca}_g\, v_{gt} \Big).
\end{align*}

\subsection*{Restricciones}

\begin{align}
\sum_{g\in\mathcal{G}} p_{gt} &\;=\; d_t
    & &\forall\, t\in\mathcal{T}, & &\text{(balance)} \\[2pt]
\mathit{Pmin}_g\, u_{gt} \;\le\; p_{gt} &\;\le\; \mathit{Pmax}_g\, u_{gt}
    & &\forall\, g\in\mathcal{G},\ t\in\mathcal{T}, & &\text{(activación)} \\[2pt]
v_{gt} &\;\ge\; u_{gt} - u_{g,t-1}
    & &\forall\, g\in\mathcal{G},\ t\in\mathcal{T}, & &\text{(arranque)} \\[2pt]
\sum_{g\in\mathcal{G}} \mathit{Pmax}_g\, u_{gt} &\;\ge\; d_t + \mathit{rr}_t
    & &\forall\, t\in\mathcal{T}, & &\text{(reserva)} \\[2pt]
-\,R_g \;\le\; p_{gt} - p_{g,t-1} &\;\le\; R_g
    & &\forall\, g\in\mathcal{G},\ t\in\mathcal{T}, & &\text{(rampa)} \\[2pt]
u_{gt} &\;\in\; \{0,1\}
    & &\forall\, g\in\mathcal{G},\ t\in\mathcal{T}, & & \notag \\[-2pt]
v_{gt} &\;\in\; \{0,1\},\quad p_{gt} \ge 0
    & &\forall\, g\in\mathcal{G},\ t\in\mathcal{T}. & & \notag
\end{align}

\noindent Para el primer periodo, $u_{g,t-1}$ y $p_{g,t-1}$ se reemplazan por los datos iniciales $u^0_g$ y $p^0_g$ en las restricciones de arranque y de rampa.

\subsection*{Comentario de modelamiento}

El corazón de esta formulación es la manera en que las variables binarias \emph{gobiernan} a la variable continua de potencia. La restricción de \textbf{activación} es en realidad una desigualdad doble que actúa como un interruptor: cuando $u_{gt}=0$ (central apagada), la cadena $\mathit{Pmin}_g\cdot 0 \le p_{gt} \le \mathit{Pmax}_g\cdot 0$ colapsa a $0 \le p_{gt} \le 0$, forzando $p_{gt}=0$; y cuando $u_{gt}=1$ (encendida), se abre la ventana física real $\mathit{Pmin}_g \le p_{gt} \le \mathit{Pmax}_g$, obligando a la central a entregar al menos su mínimo técnico y a no superar su máximo. Aquí el rol de \textbf{big-M} lo cumple naturalmente $\mathit{Pmax}_g$: es la cota superior que multiplica al binario para ``liberar'' o ``anular'' la potencia. Esta es precisamente la técnica de \emph{restricción de activación} (o restricción indicadora) vista en clase, y la razón por la que el problema no puede resolverse solo con variables continuas: sin el binario $u_{gt}$ no habría forma de imponer el mínimo técnico \emph{condicional} (``si genera, que sea al menos $\mathit{Pmin}_g$; si no, cero'').

La segunda pieza clave es el modelamiento del \textbf{arranque}, que liga periodos consecutivos. La lógica $v_{gt} \ge u_{gt} - u_{g,t-1}$ detecta la transición apagado$\to$encendido: si la central estaba apagada en $t-1$ ($u_{g,t-1}=0$) y se enciende en $t$ ($u_{gt}=1$), el lado derecho vale $1$ y obliga a $v_{gt}=1$, gatillando el pago del costo de arranque $\mathit{ca}_g$; en cualquier otro caso (seguía encendida, seguía apagada, o se apagó) el lado derecho es $\le 0$ y, como el objetivo minimiza y $v_{gt}$ tiene costo positivo, el solver la deja en $0$. No hace falta forzar la igualdad: la propia función objetivo empuja $v_{gt}$ a su mínimo. Gracias a esta restricción de acoplamiento temporal, el modelo internaliza el dilema económico central: apagar una unidad en un valle de demanda solo conviene si el ahorro en costo fijo y variable supera el costo de volver a arrancarla más tarde.

Las tres restricciones restantes traducen exigencias de negocio y de ingeniería. El \textbf{balance} garantiza que la generación total iguale exactamente la demanda en cada periodo (podría relajarse a $\ge$ si se admite sobregeneración vertida). La \textbf{reserva} exige que la capacidad máxima \emph{de las unidades encendidas} —no la potencia usada— supere la demanda más un colchón $\mathit{rr}_t$, de modo que ante una falla imprevista el sistema tenga margen para subir carga sin encender nada nuevo; es una restricción de confiabilidad que se apoya de nuevo en el binario $u_{gt}$. Finalmente, la \textbf{rampa} acota la variación de potencia entre periodos consecutivos, reflejando que las turbinas tienen inercia térmica y mecánica y no pueden cambiar su nivel de generación bruscamente. Como sugerencia de enriquecimiento, se podría añadir una restricción de \emph{tiempo mínimo encendida/apagada} (una vez que una central arranca debe permanecer operando al menos $L_g$ periodos), lo que reforzaría aún más el acoplamiento temporal entre binarios.

\fi

\end{document}
